%%%%

%%%   Самарский государственный технический университет

%%%   Кафедра прикладной математики и информатики

%%%

%%%   Преамбула для статьи в журнал "Вестник Самарского государственного

%%%   технического университета. Серия: «Физико-математические науки»"

%%%   Ориентирована под MikTEx 2.4 for Win32

%%%

%%%   Хотя сам журнал собирается под GNU/Linux на дистрибутиве TeXLive



\documentclass[11pt,a4paper,twoside]{article}



\usepackage[cp1251]{inputenc}

\usepackage[T2A]{fontenc}

\usepackage[english,russian]{babel}

\usepackage{samgtubib}

\usepackage[dvips]{graphicx}

\usepackage[multiple]{footmisc}



\usepackage{samgtu}

\renewcommand{\VestnikYear}{2009} % Год издания Вестника

\renewcommand{\VestnikNumber}{2\,(19)} % Номер Вестника

\renewcommand{\VestnikMonth}{Ноябрь} % Месяц выхода Вестника

\renewcommand{\VestnikData}{01.11.09} % Дата подписания Вестника в печать

\renewcommand{\VestnikUPL}{26,64} % Усл. печ. л.

\renewcommand{\VestnikUIL}{25,73} % Уч.-изд. л.

\renewcommand{\VestnikRegNumber}{479/09} % Регистрационный номер

\renewcommand{\VestnikOrder}{994} % Номер заказа







%

%\usepackage{qrcode}

%

%

%

%%\renewcommand{\firstpage}{99}

%%\renewcommand{\lastpage}{106}

%%\renewcommand{\thesubsection}{\arabic{subsection}.}

%%\renewcommand{\thesubsubsection}{\roman{subsection}.}

%%\newcommand{\eq}[1]{(\ref{#1})}

\newcommand{\bra}[1]{\langle{#1}\rvert}

\newcommand{\ket}[1]{\lvert{#1}\rangle}

%%\newcommand{\be}{\begin{equation}}

%%\newcommand{\end{equation}}{\end{equation}}

%%\newcommand{\mathop{\rm Tr}\nolimits}{\text{Tr}}

%%\newcommand{\ba}{\begin{eqnarray}}

%%\newcommand{\ea}{\end{eqnarray}}

%%\newcommand{\mbf}[1]{\mathbf{#1}}

%%\newcommand{\mathcal}[1]{\mathcal{#1}}

%%\newcommand{}[1]{\mathbb{#1}}

%%\newcommand{\mfr}[1]{\mathfrak{#1}}

%

%%

%%\newcommand{\totop}[1]{\raisebox{6pt}[1pt][2pt]{#1}} 

%%\newcommand{\tottop}[1]{\raisebox{12pt}[1pt][2pt]{#1}} 

%%\newcommand{\totttop}[1]{\raisebox{18pt}[1pt][2pt]{#1}} 

%%\newcommand{\tottttop}[1]{\raisebox{24pt}[1pt][2pt]{#1}} 

%%\newcommand{\totttttop}[1]{\raisebox{30pt}[1pt][2pt]{#1}} 

%\newcommand{\tobot}[1]{\raisebox{-6pt}[1pt][2pt]{#1}} 

%%\newcommand{\tobbot}[1]{\raisebox{-12pt}[1pt][2pt]{#1}} 

%%\newcommand{\tobbbot}[1]{\raisebox{-18pt}[1pt][2pt]{#1}}

%

%\graphicspath{{./00PIC/}}

%

%%\samgtupubl % для генерации pdf, отдаваемого в типографию СамГТУ

%\pdfcrop % обрезка pdf для размещения на math-net.ru

%\draft % На корректуре

%\filllastpage % Последняя страница не заполнена не полностью

%%\briefcomm % Статья является кратким сообщением



\vsgtu{1462}



\FirstPage{33}

\LastPage{42}

%%

%%

%\begin{document}





\renewcommand{\baselinestretch}{0.9}



%\hfill \begin{pspicture}(1in,1in)

%\psbarcode{http://dx.doi.org/10.14498/vsgtu1456}{}{qrcode}

%\end{pspicture}

%\vspace{-1in}

%

%\vspace{.18in}



\hfill \qrcode[hyperlink,height=.8in]{http://dx.doi.org/10.14498/vsgtu1462}

\vspace{-.8in}



%%%%%%%%%%%%%%%%%%%%%%%%%%%%%%%%%%%%%%%%%%%%%%%%%%%%%%%%%%%%%%%%%%%%%%%%%%%%%%%%%%%%

%Информация о статье и об авторах на русском и английском языках





% Номер УДК

\udc{530.145.1}

\msc{81P50}% Коды MSC см. http://www.ams.org/msc/





\rutitle{Об определении чистых состояний методом гомодинного детектирования}% Полный заголовок

{Об определении чистых состояний методом гомодинного детектирования}



\entitle{On the determination of pure quantum states by~the~homodyne detection}



\ruauthor{А.~И.~Днестрян}{Д\,н\,е\,с\,т\,р\,я\,н~А.~И.}

\enauthor{A.~I.~Dnestryan}{D\,n\,e\,s\,t\,r\,y\,a\,n~A.~I.}



\ruaffil{\rumfti.}

\enaffil{\enmfti.}







\ruauthordetails{1}{

\fullauthor{\rupid{73344}{Андрей Игоревич Днестрян}} \regalia{\mem{dnestor@inbox.ru}} \position{аспирант} \dept{каф. высшей математики} }





\enauthordetails{1}{

\fullauthor{\enpid{73344}{Andrey I. Dnestryan}} \regalia{\mem{dnestor@inbox.ru}} \position{Postgraduate Student}

\dept{Dept. of Higher Mathematics}}







\rukeywords{квантовая томограмма, квантовое состояние, оператор плотности, волновая функция}



\ruabstract{В работе обсуждаются методы реконструкции волновой  функции чистого состояния квантовой системы по известной оптической томограмме состояния.

Оптическая квантовая томограмма представляет собой однопараметрическое распределение вероятностей с параметром $\theta$.

Волновая функция чистого состояния выражается через оптическую томограмму, если последняя известна для любых значений

$\theta$. Однако оптическая томограмма определяется из эксперимента гомодинного детектирования, где $\theta$ фиксированно.

Поэтому оптическая томограмма может быть известна лишь для нескольких дискретных значений параметра. 

Мы приводим  приближенные методы определения волновой функции квантового состояния по неполной информации о его томограмме, представляющие собой  развитие уже существующих методов.

}



\enabstract{The methods of reconstruction of the wave function of a pure state of a quantum system by quadrature distribution measured experimentally by the homodyne detection are considered. Such distribution is called optical tomogram of a state and containes one parameter $\theta$. Wave function of a state is determined exactly by its optical tomogram if last one is known for all $\theta$. But one can obtain optical tomogram from experiment of homodyne detection only for discrete number of $\theta$. We introduce some approximate methods of reconstructing the state by such information about its optical tomogram.}



\enkeywords{quantum tomography, quantum state, density operator, wave function}





\date{21/XI/2015} % Дата поступления статьи в редакцию.

\revisiondate{24/II/2016} % В окончательном виде

\accepteddate{26/II/2016}







\makerutitle







\Section{Введение}

В конце XX века было предложено новое вероятностное представление квантовой механики, в которой квантовые состояния связываются со стандартными плотностями распределений вероятностей "--- так называемыми томографическими плотностями или квадратурными распределениями. 

Этот формализм, содержащий такую же информацию, как волновая функция и матрица плотности, основан на томографическом подходе к~измерению квантовых состояний \cite{bert,raymer}. 

Суть его в~следующем: каждому состоянию квантовой системы с матрицей плотности $\hat\rho$ ставится в соответствие плотность распределения наблюдаемой "--- квадратурной компоненты $\hat X=\mu\hat q +\nu\hat p$:

\begin{equation}\label{q}

w(X,\mu,\nu)=\mathop{\rm Tr}\nolimits{\hat\rho\delta(X-\mu\hat q-\nu\hat p)},

\end{equation}

где $\hat q$ и $\hat p$ есть обычные операторы координаты и импульса. 

Функция \eqref{q} называется симплектической квантовой томограммой состояния $\hat\rho$. 

Частным случаем симплектической томограммы является экспериментально измеримая \cite{raymer} оптическая томограмма

$$w(X,\theta)=w(X,\mu=\cos\theta,\nu=\sin\theta).$$







В работе \cite{manko} было установлено, что для чистого состояния с координатной волновой функцией $\psi(x)$ его симплектическая томограмма выражается следующим образом:

\begin{equation}\label{ma}

w(X,\mu,\nu)=\big|\hat F_{\mu,\nu}\; [\psi]\big|^2(X),

\end{equation}

где $\hat F_{\mu,\nu}$ есть линейный интегральный оператор в $\mathcal{L}^2(\mathbb{R})$, подробно изученный в \cite{dn}. В данной статье мы рассмотрим вопрос обратимости отображения \eqref{ma} в случае чистых состояний.



\smallskip



\Section{Реконструкция чистого состояния по его квадратурному распределению}

Результатом всякой процедуры измерений над квантовой

системой может быть только распределение вероятностей. Поскольку

квантовое состояние содержит всю доступную информацию о квантовой системе, мы безусловно можем, отталкиваясь от этого состояния,

рассчитать все распределения вероятностей \cite{slyah}. Зададим обратный вопрос:

возможно ли использовать набор вероятностных распределений для

реконструкции квантового состояния?



Этот вопрос возвращает нас к раннему периоду развития квантовой механики, в частности, к обзорной статье В.~Паули \cite{pauli1}. 

Он интересовался

вопросом, можно ли найти амплитуду и фазу волновой функции, зная

вероятности распределений по координате и импульсу. Паули не дал

ответа на этот вопрос. Однако простые контрпримеры 

(см., например, [\citen{reichenbach}--\citen{welsch}]) показывают, что

в общем случае это невозможно. 

На самом деле нужно знать больше распределений, чем эти два.



Как известно, функция Вигнера содержит всю информацию о квантовом состоянии,

 поэтому мы можем восстановить все квадратурные распределения с помощью преобразования Радона.



Преобразование вида

\begin{equation}\label{radon}

w(x,\mu,\nu)= { \int_{-\infty}^{\infty} \int_{-\infty}^{\infty}} W(q,p)\delta(x-\mu q-\nu p)\;dqdp

\end{equation}

называется преобразованием Радона \cite{radon} функции $W(q,p)$. 

Функция

 \eqref{radon} совпадает с симплектической томограммой \cite{manko}.

 Преобразование Радона обратимо, обращение преобразования \eqref{radon} выглядит следующим образом:

$$

W(q,p)=\frac{1}{(2\pi)^2} {\int_{-\infty}^{\infty} \int_{-\infty}^{\infty} \int_{-\infty}^{\infty}} w(x,\mu,\nu)e^{i(x-\mu q-\nu p)}\;dxd\mu d\nu.

$$

Обратимость преобразования Радона здесь позволяет выразить функцию Вигнера  $W(q,p)$ через симплектическую томограмму состояния $w(x,\mu,\nu)$ при условии, что последняя известна для всех действительных $\mu,\nu$. 

Симплектическая томограмма обладает замечательным свойством однородности

$$w(\lambda X,\lambda\mu,\lambda\nu)=\frac{1}{|\lambda|}w(X,\mu,\nu),

$$

 используя которое, можно восстановить функцию Вигнера, зная лишь оптическую 

 томограмму $w(x,\theta)$ для любого $\theta \in [0, \pi]$ \cite{dariano,leonhardt2}. 

 В~этом случае функция Вигнера выражается по формуле

\begin{equation}

W(q,p)=\frac{1}{(2\pi)^2}\int_{0}^{\pi}d\theta

\int_{-\infty}^{+\infty}|t|\;dt\int_{-\infty}^{+\infty}w(x,\theta)\exp{it(x-q\cos\theta -p\sin\theta )}\;dx.

\end{equation}



Однако реально невозможно измерить оптическую томограмму (квадратурное распределение)

 $w(x,\theta)$ для {\it{любых}} значений фазы $\theta$. Возможно лишь провести гомодинное

 детектирование для различных, но фиксированных значений фазы $\theta$. 

 На практике, когда такой подход реализуется, получается ансамбль распределений 

 $\{w(x,\theta_1),\,w(x,\theta_2),\dots, w(x,\theta_{N})\}$, который с некоторой 

 погрешностью можно считать за истинный непрерывный ансамбль. 

 Затем по этому

 ансамблю численно с помощью обратного преобразования Радона восстанавливается 

 функция Вигнера исследуемого состояния.

 

В эксперименте фазу $\theta$ изменяют, меняя разность хода поступающих на 

светоделитель лучей. Удобно выбирать значения фаз эквидистантными $\theta_n={\pi (n-1)}/{N}$.

 Поэтому для осуществления обратного преобразования Радона необходимо, чтобы $n$ пробегало

 значения от 1 до $N$. Ясно, что точность вычислений увеличивается с ростом $N$, т. к.

 численное интегрирование представляет собой суммирование с шагом $h \propto N^{-1}$.

  Однако зачастую проведение экспериментов со значениями $N\geqslant 10$ бывает затруднительным.







Вместе с этим количество фаз $N$ еще и качественно влияет на результат. 

В~работе \cite{leonhardt1}

было показано, что для точной реконструкции состояния в оптической гомодинной

 томографии размерность матрицы плотности в представлении Фока должна быть равна количеству фаз $\theta_n$, для которых было произведено гомодинное детектирование. Также был получен простой способ оценки ошибок, если фактическая размерность матрицы плотности больше, чем число фаз, используемых в эксперименте.



Еще один метод реконструкции волновой функции состояния был дан в \cite{orlow}.

 Он основан на интегральных представлениях коэффициентов в разложении волновой функции. Развивая идею, высказанную авторами, мы предлагаем следующий метод.



Пусть произведен эксперимент гомодинного детектирования некого состояния. 

На выходе эксперимента для фиксированного значения фазы $\theta$ мы имеем относительно большой ($\sim 10^5$) набор точек $x_j$ "--- измеренных квадратурных компонент, по которому строится гистограмма. 

Эта гистограмма определяет распределение $w(x,\theta)$. Предположим, что чистое состояние $$\hat\rho=\ket{\psi}\bra{\psi},$$ подаваемое на гомодинный детектор, представимо конечной суммой фоковских состояний, т.~е. волновая функция может быть представлена в виде

\begin{equation}\label{fok}

\ket{\psi}=\sum\limits_{n=0}^{N}c_n\ket{n},

\end{equation}

где

\begin{equation}\label{dva1}

\langle x\ket{n}=\frac{1}{\pi^{1/4}\sqrt{2^n n!}}H_n(x)e^{-x^2/2}.

\end{equation}

Фоковские состояния образуют ортонормированный базис

$$\langle{n}\ket{m}=\delta_{nm},$$ поэтому коэффициенты $c_n$ удовлетворяют

$$\sum_{n=0}^{N}|c_n|^2=1.$$

Обозначая $\langle x\ket{n}=\psi_n(x)$ и подставляя разложение \eqref{fok} в 

равенство \eqref{ma}, мы получаем оптическую томограмму суперпозиции фоковских состояний \cite{dn}

\begin{multline}

w(x, {\theta})=\big|\hat F_{\cos\theta,\sin\theta}\; [\psi]\big|^2 = \bigg|\sum\limits_{n=0}^{N}c_n  \psi_n (x)  e^{-in{\theta}}\bigg|^2 =

\\ 

= \sum\limits_{n,m \geqslant 0}^N\text{Re} \big(c_n c_m^*e^{i{\theta}(m-n)}\big) \psi_n (x) \psi_m (x) = 

\\

= \sum\limits_{n=0}^{N}|c_n|^2  \psi_n^2 (x)  + 2 \sum\limits_{0\leqslant m < l\leqslant N }\text{Re} \big(c_m c_l^* e^{i{\theta}(l-m)}\big)\psi_m (x)\psi_l (x)=

\\

\label{dva}

=\sum\limits_{n,m\geqslant 0}^Nc_n c_m^* e^{i{\theta}(m-n)}\psi_n (x)\psi_m (x)

\end{multline}

или через амплитуду и фазу коэффициентов $c_n=|c_n|e^{i\varphi_n}$:

\begin{equation}\label{dvad}

w(x, {\theta})=\sum\limits_{0\leqslant n,m\leqslant N}|c_n|| c_m| \cos\left(\theta(m-n)+(\varphi_n-\varphi_m)\right)\psi_n (x)\psi_m (x).

\end{equation}

Произведение

 волновых функций фоковских состояний $\psi_n (x) \psi_m (x)$ раскладывается в сумму \cite{orlow}

\begin{equation}

\label{dvad5}

\psi_n (x) \psi_m (x)=2^{1/4}\sum\limits_{k=0}^{n+m}\beta_k^{n,m}\psi_k (\sqrt{2}x),

\end{equation}

где коэффициенты $\beta_k^{n,m}$ не равны нулю, только если числа $n+m$ и $k$ 

одной четности, при этом для незануляющихся коэффициентов

$$

\beta_k^{n,m}=\left(\frac{2}{\pi}\right)^{1/4}\sqrt\frac{n!m!}{k!}2^{-2q-(k+1)/2}(-1)^q\sum\limits_{j=0}^{\text{min}(n,m,q)}\frac{(-4)^j(n+m-2j)!}{j!(n-j)!(m-j)!(q-j)!},$$

где $q=(n+m-k)/2$. В \hyperlink{tabl:new}{таблице} представлены коэффициенты $\beta_k^{n,m}$ 

для $1\leqslant n \leqslant 2$, $1\leqslant m \leqslant 5$, определенные численными методами по формуле

\begin{equation}

\label{new-tab}

\beta_k^{n,m}=2^{1/4}\int _{-\infty}^{+\infty}\psi_n (x) \psi_m (x)\psi_k (\sqrt{2}x)dx,

\end{equation}

следующей напрямую из \eqref{dvad5}.

Отсюда следует, что \eqref{dvad}  можно переписать так:

\begin{equation}

\label{dvad1}

w(x, {\theta})=2^{1/4}\!\!\!\!\!\!\sum\limits_{0\leqslant n,m\leqslant N}\sum\limits_{k=0}^{n+m}|c_n|| c_m| \cos\left(\theta(m-n)+(\varphi_n-\varphi_m)\right)\beta_k^{n,m}\psi_k (\sqrt{2}x).\end{equation}

Поскольку $\psi_k(\sqrt{2}x)$ образуют ортогональный базис в $\mathcal{L}^2(\mathbb{R})$,  равенство \eqref{dvad1} представляет из себя разложение $w(x,\theta)$  в ряд Фурье, а коэффициенты этого ряда находятся из условия

\begin{multline}\label{dvad2}

\sum\limits_{0\leqslant n,m\leqslant N}|c_n|| c_m| \cos\left(\theta(m-n)+(\varphi_n-\varphi_m)\right)\beta_k^{n,m}=\\ =

2^{1/4}\left(w(x,\theta),\psi_k (\sqrt{2}x)\right).

\end{multline}

Скалярное произведение в правой части \eqref{dvad2} определено  в $\mathcal{L}^2(\mathbb{R})$, т.~е.

$$\left(w(x,\theta),\psi_k (\sqrt{2}x)\right)=\int_{-\infty}^{+\infty}w(x,\theta)\psi_k (\sqrt{2}x)dx.

$$

Индекс $k$ в \eqref{dvad2} пробегает $2N+1$ значений:  $0\leqslant k \leqslant 2N$. 

Тем самым для нахождения коэффициентов $c_n$ разложения \eqref{fok} мы имеем систему из $2N+1$ уравнений, повторяющих \eqref{dvad2} для разных значений $k$:

\begin{equation}

 \label{dvad3}

          \left\{

           \begin{array}{l}

           \sum\limits_{0\leqslant n,m\leqslant N}|c_n|| c_m| \cos\left(\theta(m-n)+(\varphi_n-\varphi_m)\right)\beta_{2N}^{n,m}=

           \\

\hspace{6cm} =           

2^{1/4}\left(w(x,\theta),\psi_{2N} (\sqrt{2}x)\right);\\[3mm]

            \sum\limits_{0\leqslant n,m\leqslant N}|c_n|| c_m| \cos\left(\theta(m-n)+(\varphi_n-\varphi_m)\right)\beta_{2N-1}^{n,m}=

                       \\

\hspace{6cm} =           

2^{1/4}\left(w(x,\theta),\psi_{2N-1} (\sqrt{2}x)\right);\\

\hspace{3cm}            \dots\ldots\ldots\ldots \\[4mm]

           \sum\limits_{0\leqslant n,m\leqslant N}|c_n|| c_m| \cos\left(\theta(m-n)+(\varphi_n-\varphi_m)\right)\beta_k^{n,m}=

                                  \\

\hspace{6cm} =           

2^{1/4}\left(w(x,\theta),\psi_{k} (\sqrt{2}x)\right);\\[3mm]

\hspace{3cm}  \dots\ldots\ldots\ldots \\[4mm]         

 \sum\limits_{0\leqslant n,m\leqslant N}|c_n|| c_m| \cos\left(\theta(m-n)+(\varphi_n-\varphi_m)\right)\beta_0^{n,m}=

                        \\

\hspace{6cm} =           

2^{1/4}\left(w(x,\theta),\psi_{0} (\sqrt{2}x)\right).\\

           \end{array}

\right.

  \end{equation}





\begin{table}[h]

\begin{center}

\begin{minipage}{.8\textwidth}

\hypertarget{tabl:new}{}

 \small

 {\bf Таблица коэффициентов $\boldsymbol{\beta_k^{n,m}}$, вычисленных по \eqref{new-tab}. 

Верхняя таблица содержит коэффициенты $\boldsymbol{\beta_k^{1,m}}$, т. е. при $\boldsymbol{n=1}$,

 нижняя таблица "--- при $\boldsymbol{n=2}$ [The table contains the coefficients $\boldsymbol{\beta_k^{n,m}}$ calculated by the Eq.~\eqref{new-tab}.

The upper part of the table contains the coefficients $\boldsymbol{\beta_k^{1,m}}$; 

\centerline{The lower part of the table contains the coefficients $\boldsymbol{\beta_k^{2,m}}$]} }

 

\centering

\begin{tabular}{|c||c|c|c|c|c|c|}

\hline

$n=1$ & $m=1$ & $m=2$ & $m=3$ & $m=4$ & $m=5$  \\

\hline

\hline

$k=1$ & $0$ & $0.1579$ & $0$ & $-0.2051$ & $0$\\

\hline

$k=2$ & $0.4466$ & $0$ & $0$ & $0$ & $-0.1528$\\

\hline

$k=3$ & $0$ & $0.3868$ & $0$ & $-0.1116$ & $0$\\

\hline

$k=4$ & $0$ & $0$ & $0.3158$ & $0$ & $-0.1765$\\

\hline

$k=5$ & $0$ & $0$ & $0$ & $0.2496$ & $0$\\

\hline

$k=6$ & $0$ & $0$ & $0$ & $0$ & $0.1934$\\

\hline

%\end{tabular}

%\end{center}

%\end{table} \\

%\begin{table}[b]

%\begin{center}

%\begin{tabular}{|c||c|c|c|c|c|c|}

\hline

$n=2$ & $m=1$ & $m=2$ & $m=3$ & $m=4$ & $m=5$  \\

\hline

\hline

$k=1$ & $0.1579$ & $0$ & $0.0967$ & $0$ & $-0.1621$\\

\hline

$k=2$ & $0$ & $0.2233$ & $0$ & $-0.0483$ & $0$\\

\hline

$k=3$ & $0.3867$ & $0$ & $0.1579$ & $0$ & $-0.1324$\\

\hline

$k=4$ & $0$ & $0.3867$ & $0$ & $-0.0558$ & $0$\\

\hline

$k=5$ & $0$ & $0$ & $0.3531$ & $0$ & $-0.0395$\\

\hline

$k=6$ & $0$ & $0$ & $0$ & $0.3058$ & $0$\\

\hline

$k=7$ & $0$ & $0$ & $0$ & $0$ & $0.2558$\\

\hline

\end{tabular}

\end{minipage}

\end{center}

\end{table}





Система уравнений \eqref{dvad3}, вообще говоря, содержит $2N+2$ неизвестных $\bigl(\{c_k\}_{k=0}^{N}$ и $\{\varphi_k\}_{k=0}^{N}\bigr)$ и $2N+1$ уравнений, 

 однако она не является недоопределенной. 

 Дело в том, что искомая 

 волновая функция $\psi(x)$ в любом случае может быть определена с

 точностью до постоянной фазы. 

 По этой причине мы фиксируем (зануляем)   фазу, например, коэффициента $c_N$, а все остальные фазы считаем  относительно фазы $\varphi_N=0$. 

 После этого предположения число  неизвестных системы \eqref{dvad3} становится равным $2N+1$.

 



 

На первый взгляд, решение системы уравнений \eqref{dvad3} представляется 

очень трудным процессом, однако есть один момент, заметно упрощающий 

ее решение. Суть в следующем: коэффиценты $\beta_k^{n,m}=0$ при $n+m <k$. 

Следовательно, первые уравнения перепишутся заметно проще:

\begin{equation} \label{dvad4}

          \left\{

           \begin{array}{l}

           c_N^2\beta_{2N}^{N,N}=

2^{1/4}\left(w(x,\theta),\psi_{2N} (\sqrt{2}x)\right);\\[3mm]

           2c_N c_{N-1} \cos\left(\theta-\varphi_{N-1}\right)\beta_{2N-1}^{N,N-1}=

2^{1/4}\left(w(x,\theta),\psi_{2N-1} (\sqrt{2}x)\right);\\

\hspace{3cm}\ldots\ldots\ldots\ldots \\

           \end{array}

\right.

  \end{equation}

Отсюда видно, что $c_N$ находится из первого уравнения, а $c_{N-1}$ 

выражается через $c_N$ и т.~д. 

Таким образом, данная методика позволяет  определить состояние $\psi(x)$ в представлении \eqref{fok}, т.~е. в виде суперпозиции первых $N$ фоковских состояний.

Отметим, что исходный метод \cite{orlow} решает данную задачу при $N=2$.







Перейдем к следующему методу реконструкции волновой функции, который является развитием метода, представленного в \cite{dnn}.

 Вводя коэффициент $$\alpha_n=\frac{1}{\sqrt[4]{\pi}\sqrt{2^n n!}},

 $$ перепишем равенство \eqref{dva}, используя \eqref{dva1}:

\begin{multline}

w(x,{\theta})= \sum\limits_{n=0}^{N}\alpha_n^2|c_n|^2  H_n^2 (x)e^{-x^2}  + \\

+  2 \sum\limits_{0\leqslant m < l \leqslant N}\alpha_m \alpha_l \text{Re} \big(c_m c_l^* e^{i{\theta}(l-m)}\big)H_m (x)H_l (x)e^{-x^2}.

\end{multline}

Домножим обе части этого равенства на $e^{x^2}$ и получим

\begin{equation}\label{dva2}

e^{x^2}w(x,\theta)= \sum\limits_{n=0}^{N}\alpha_n^2|c_n|^2  H_n^2 (x)  +  2 \sum\limits_{0\leqslant m < l \leqslant N}\alpha_m \alpha_l\text{Re} \big(c_m c_l^* e^{i{\theta}

(l-m)}\big)H_m (x)H_l (x).

\end{equation}

В этом равенстве справа стоит многочлен степени $2N$, его коэффициенты 

можно определить, дифференцируя этот многочлен пошагово:

$$

           \left\{

           \begin{array}{rcl}

\displaystyle \frac{d^{2N}}{dx^{2N}}\big(e^{x^2}w(x,{\theta})\big)&=& \alpha_N^2 (2N)! 2^{2N} |c_N|^2;\\[3mm]

\displaystyle \frac{d^{2N-1}}{dx^{2N-1}}\big(e^{x^2}w(x,{\theta})\big)&=& \alpha_N^2 (2N)! 2^{2N} |c_N|^2 x +{}

\\

& & {}+ 2\alpha_N^2 N(N-1) 2^{N-2 }2^N (2N-1)!|c_N|^2 + {} \\

            &  &{}+ 2^N 2^{N-1}(2N-1)!\alpha_N \alpha_{N-1}\text{Re}\left(c_N c_{N-1}^* e^{i{\theta}}\right);\\

& &            \ldots \\[2mm]

           e^{x^2}w(x,\theta)&=& \displaystyle \sum\limits_{n=0}^{N}\alpha_n^2|c_n|^2  H_n^2 (x) +{} \\

           & & {} + \displaystyle 2 \sum\limits_{0\leqslant m < l \leqslant N}\alpha_m \alpha_l \text{Re} \big(c_m c_l^* e^{i{\theta}(l-m)}\big)H_m (x)H_l (x).

           \end{array}

\right.

  $$

Система  содержит $2N+1$ уравнений и $2N+2$ неизвестных, доопределим эту 

систему условием нормировки волновой функции

$$\sum\limits_{n=0}^{N}|c_n|^2 = 1.$$

Тогда количество неизвестных и уравнений совпадает и равно $2N+2$.

Решение системы дает коэффициенты $c_n$ в разложении \eqref{fok}.



\smallskip



\Section{Заключение}

В данной статье развиты уже существующие методы реконструкции чистого квантового состояния 

по неполной информации о его квантовой томограмме. 

Состояние описывается волновой  функцией, которая в рамках данной работы аппроксимируется суммой из $N$ фоковских  состояний. 

 Отметим, что по сравнению с начальными методами, где $N$ ограничено, в данной работе это число произвольно. 

С~ростом количества слагаемых данной суммы в силу сходимости можно добиться сколь угодно  большой точности приближений волновой функции. Также метод, развитый в данной работе, имеет преимущество перед обратным преобразованием

Радона, так как последнее осуществимо только если известны квадратурные распределения $w(x,\theta)$ на сетке $\theta_n$,

покрывающей отрезок $[0;\pi]$, тогда как метод, предложенный в работе, требует, чтобы было известно квадратурное распределение

 лишь для одного значения $\theta$. Это преимущество существенно при проведении экспериментов гомодинного детектирования. 



Отметим, что на практике всегда может быть известна лишь неполная информация о томограммме, 

поскольку томограмма является экспериментально наблюдаемой величиной, и результат эксперимента 

представляет собой набор значений томограммы в разных точках. 

Этот факт обуславливает 

актуальность задачи, рассмотренной в данной работе.





\thanks{{\bf ORCID}



{\bf Андрей Игоревич Днестрян:} \url{http://orcid.org/0000-0002-9381-2133}



}



\newpage



\RusRef



\begin{thebibliography}{99}



\Bibitem{bert}

\by Vogel~K., Risken~H.

\paper Determination of quasiprobability distributions in  terms of probability distributions for the rotated  quadrature phase

\jour Phys. Rev. A

\yr 1989

\issue 5

\vol 40

\pages 2847--2855

\crossref{10.1103/physreva.40.2847}



\Bibitem{raymer}

\by  Smithey~D.~T., Beck~M., Raymer~M.~G., Faridani~A.

\paper Measurement of the Wigner distribution and the density matrix of a light mode using optical homodyne tomography: Application to squeezed states and the vacuum

\jour Phys. Rev. Lett.

\yr 1993

\vol 70

\issue 9 

\pages 1244--1247

\crossref{10.1103/physrevlett.70.1244} 







\Bibitem{manko}

\by   Mancini~S.,  Man'ko~V.~I.,  Tombesi~P.

\paper Symplectic tomography as classical approach to quantum systems

 \jour Phys. Lett. A

\yr 1996

\vol 213

\issue 1--2

\pages 1--6

\crossref{10.1016/0375-9601(96)00107-7}



\RBibitem{dn}

\by  Амосов~Г.~Г., Днестрян~А.~И.

\paper О спектре семейства  интегральных операторов, определяющих квантовую  томограмму

\jour Труды МФТИ

\yr 2011

\vol 3

\issue 1

\pages 5--9

\elib{http://elibrary.ru/item.asp?id=16904483}







\Bibitem{slyah}

\by  Schleich~W.~P.

\book Quantum Optics in Phase Space

\publ Verlag 

\publaddr Berlin

\yr 2001,

%\totalpages 

xx+695~pp

\crossref{10.1002/3527602976}





\Bibitem{pauli1}

\by Pauli~W.

\paper Die allgemeinen Prinzipien der Wellenmechanik

\zmath{https://zbmath.org/?q=an:0007.13504}

\jour Handbuch Physik 

\vol 24

\volinfo Tl.~1

\yr 1933

\pages 83--272;

%\by 

 Pauli~W.

%\paper 

Die allgemeinen Prinzipien der Wellenmechanik

%\inbook 

/

{\it Die allgemeinen Prinzipien der Wellenmechanik}

%\serial 

/ 

Neu herausgegeben und mit historischen Anmerkungen versehen von Norbert Straumann;

ed. Professor Dr. Norbert Straumann.

Berlin:

Springer, 1990.

 pp.~21--192

\crossref{10.1007/978-3-642-61287-9_2} 

%\moreref

%\by  Pauli~W.

%\paper Die allgemeinen Prinzipien der Wellenmechanik

%\inbook Die allgemeinen Prinzipien der Wellenmechanik

%\ed Professor Dr. Norbert Straumann

%\serial Neu herausgegeben und mit historischen Anmerkungen versehen von Norbert Straumann

%\pages 21--192

%\yr 1990

%\publ Springer 

%\publaddr Berlin 

%\crossref{10.1007/978-3-642-61287-9_2}





\Bibitem{reichenbach}

\by  Reichenbach~H.

\book Philosophic Foundations of Quantum Mechanics

\yr 1944

\publ University of California Press

\publaddr Berkeley and Los Angeles 

\totalpages x+182



\Bibitem{vogt}

\by Vogt~A.

\paper Position and Momentum Distributions do not Determine the Quantum Mechanical State

\inbook Mathematical Foundations of Quantum Theory

\yr 1978

\publ Academic Press

\publaddr New York

\pages 365--372

\crossref{10.1016/b978-0-12-473250-6.50024-8}





\Bibitem{freyberger}

\by  Freyberger~M., Bardroff~P.~J., Leichle~C., Schrade~G.,  Schleich~W.~P.

\paper The art of measuring quantum states

 \jour Physics World

\yr 1997

\vol 10

\issue 11

\pages 41--46

\crossref{10.1088/2058-7058/10/11/31}





\Bibitem{slyahraymer}

\by  Schleich~W.~P., Raymer~M.~G.

\paper Special issue on quantum state preparation and measu\-rement

\jour J. Mod. Opt.

\yr 1997

\issue 44

\issue 11--12 

\pages 2021--2022

\crossref{10.1080/09500349708231863}



\Bibitem{leibfried}

\by  Leibfried~D., Pfau~T., Monroe~C.

\paper Shadows and Mirrors: Reconstructing Quantum States of Atom Motion

 \jour Physics Today

\yr 1998

\vol 51

\issue 4

\pages 22--28

\crossref{10.1063/1.882256}



\Bibitem{welsch}

\by Welsch~D.-G., Vogel~W.,  Opatrn\'y~T.

\paper II Homodyne Detection and Quantum-State Reconstruction

\yr 1999 

\inbook Progress in Optics 

\pages 63--211

\bookvol 39

\ed E.~Wolf

\publ North-Holland

\publaddr Amsterdam

\crossref{10.1016/S0079-6638(08)70389-5}



\Bibitem{radon}

\by  Radon~J.

\paper \"Uber die Bestimmung von Funktionen durch ihre Integralwerte l\"angs gewisser Mannigfaltigkeiten

 \jour Ber. Verh. S\"achs. Akad. Wiss. Leipzig, Math. Nat. kl.

\yr 1917

\vol 69

\pages 262--277

\elink Available at 

 \url{http://people.csail.mit.edu/bkph/courses/papers/Exact_Conebeam/Radon_Deutsch_1917.pdf}    (February 24, 2016);

%\moreref

%\by  

Radon~J. 

%\paper 

\"Uber die Bestimmung von Funktionen durch ihre Integralwerte l\"angs gewisser Mannigfaltigkeiten

%\inbook 

/

{\it Proceedings of Symposia in Applied Mathematics}.

%\bookvol 

vol.~27;

%\ed 

ed.

Lawrence~A. Shepp.

%\publaddr 

Providence:

%\publ 

Amer. Math. Soc.,

%\yr 

1983,

%\pages 

pp.~71--86

\crossref{10.1090/psapm/027/692055}





\Bibitem{dariano}

\by   D'Ariano~G.~M., Mancini~S., Man'ko~V.~I., Tombesi~P.

\paper Reconstructing the density operator by using generalized field quadratures

 \jour Quantum Semiclass. Opt. 

\yr 1996

\issue 5

\vol 8

\pages 1017--1027

\crossref{10.1088/1355-5111/8/5/007}





\Bibitem{leonhardt2}

\by Leonhardt~U.,  Paul~H., D'Ariano~G.~M.

\paper Tomographic reconstruction of the density matrix via pattern functions

 \jour Phys. Rev.~A

\yr 1995

\vol 52

\issue 6

\pages 4899--4907

\crossref{10.1103/physreva.52.4899}



\Bibitem{leonhardt1}

\by  Leonhardt~U., Munroe~M.

\paper Number of phases required to determine a quantum state in optical homodyne tomography

 \jour  Phys. Rev.~A

\yr 1996

\vol 54

\issue 4

\pages 3682--3684

\crossref{10.1103/physreva.54.3682}





\Bibitem{orlow}

\by Or{\l}owski~A.,  Paul~H.

\paper Phase retrieval in quantum mechanics

\jour Phys. Rev. A

\yr 1994

\vol 50

\issue 2

\pages R921--R924

\crossref{10.1103/physreva.50.r921}









\RBibitem{dnn}

\by Амосов~Г.~Г., Днестрян~А.~И.

\paper О восстановлении чистого состояния по неполной информации о~его оптической томограмме

\jour Изв. вузов. Матем.

\yr 2013

\issue 3

\pages 62--67

\mathnet{http://mi.mathnet.ru/ivm8784}



\end{thebibliography}



\RuDate





%

% \newpage



\smallskip

 

 

 \makeentitle

 

\thanks{{\bf ORCID}



{\bf Andrey I. Dnestryan:} \url{http://orcid.org/0000-0002-9381-2133}



}



\EngRef



\begin{thebibliography}{99}



\Bibitem{bert:eng}

\by Vogel~K., Risken~H.

\paper Determination of quasiprobability distributions in  terms of probability distributions for the rotated  quadrature phase

\jour Phys. Rev. A

\yr 1989

\issue 5

\vol 40

\pages 2847--2855

\crossref{10.1103/physreva.40.2847}



\Bibitem{raymer:eng}

\by  Smithey~D.~T., Beck~M., Raymer~M.~G., Faridani~A.

\paper Measurement of the Wigner distribution and the density matrix of a light mode using optical homodyne tomography: Application to squeezed states and the vacuum

\jour Phys. Rev. Lett.

\yr 1993

\vol 70

\issue 9 

\pages 1244--1247

\crossref{10.1103/physrevlett.70.1244} 





\Bibitem{manko:eng}

\by   Mancini~S.,  Man'ko~V.~I.,  Tombesi~P.

\paper Symplectic tomography as classical approach to quantum systems

 \jour Phys. Lett. A

\yr 1996

\vol 213

\issue 1--2

\pages 1--6

\crossref{10.1016/0375-9601(96)00107-7}



\Bibitem{dn:eng}

\lang In Russian

\by  Amosov~G.~G., Dnestryan~A.~I.

\paper On the spectrum of a family of integral operators defining the quantum tomogram

\jour Trudy MFTI

\yr 2011

\vol 3

\issue 1

\pages 5--9





\Bibitem{slyah:eng}

\by  Schleich~W.~P.

\book Quantum Optics in Phase Space

\publ Verlag 

\publaddr Berlin

\yr 2001,

%\totalpages 

xx+695~pp

\crossref{10.1002/3527602976}





\Bibitem{pauli1:eng}

\by Pauli~W.

\paper Die allgemeinen Prinzipien der Wellenmechanik

\zmath{https://zbmath.org/?q=an:0007.13504}

\jour Handbuch Physik 

\vol 24

\volinfo Tl.~1

\yr 1933

\pages 83--272;

%\by 

 Pauli~W.

%\paper 

Die allgemeinen Prinzipien der Wellenmechanik,

%\inbook 

{\it Die allgemeinen Prinzipien der Wellenmechanik},

%\serial 

 Neu herausgegeben und mit historischen Anmerkungen versehen von Norbert Straumann;

ed. Professor Dr. Norbert Straumann.

Berlin,

Springer, 1990.

 pp.~21--192

\crossref{10.1007/978-3-642-61287-9_2} 

%\moreref

%\by  Pauli~W.

%\paper Die allgemeinen Prinzipien der Wellenmechanik

%\inbook Die allgemeinen Prinzipien der Wellenmechanik

%\ed Professor Dr. Norbert Straumann

%\serial Neu herausgegeben und mit historischen Anmerkungen versehen von Norbert Straumann

%\pages 21--192

%\yr 1990

%\publ Springer 

%\publaddr Berlin 

%\crossref{10.1007/978-3-642-61287-9_2}





\Bibitem{reichenbach:eng}

\by  Reichenbach~H.

\book Philosophic Foundations of Quantum Mechanics

\yr 1944

\publ University of California Press

\publaddr Berkeley and Los Angeles 

\totalpages x+182



\Bibitem{vogt:eng}

\by Vogt~A.

\paper Position and Momentum Distributions do not Determine the Quantum Mechanical State

\inbook Mathematical Foundations of Quantum Theory

\yr 1978

\publ Academic Press

\publaddr New York

\pages 365--372

\crossref{10.1016/b978-0-12-473250-6.50024-8}





\Bibitem{freyberger:eng}

\by  Freyberger~M., Bardroff~P.~J., Leichle~C., Schrade~G.,  Schleich~W.~P.

\paper The art of measuring quantum states

 \jour Physics World

\yr 1997

\vol 10

\issue 11

\pages 41--46

\crossref{10.1088/2058-7058/10/11/31}





\Bibitem{slyahraymer:eng}

\by  Schleich~W.~P., Raymer~M.~G.

\paper Special issue on quantum state preparation and measu\-rement

\jour J. Mod. Opt.

\yr 1997

\issue 44

\issue 11--12 

\pages 2021--2022

\crossref{10.1080/09500349708231863}



\Bibitem{leibfried:emg}

\by  Leibfried~D., Pfau~T., Monroe~C.

\paper Shadows and Mirrors: Reconstructing Quantum States of Atom Motion

 \jour Physics Today

\yr 1998

\vol 51

\issue 4

\pages 22--28

\crossref{10.1063/1.882256}



\Bibitem{welsch:eng}

\by Welsch~D.-G., Vogel~W.,  Opatrn\'y~T.

\paper II Homodyne Detection and Quantum-State Reconstruction

\yr 1999 

\inbook Progress in Optics 

\pages 63--211

\bookvol 39

\ed E.~Wolf

\publ North-Holland

\publaddr Amsterdam

\crossref{10.1016/S0079-6638(08)70389-5}



\Bibitem{radon:eng}

\by  Radon~J.

\paper \"Uber die Bestimmung von Funktionen durch ihre Integralwerte l\"angs gewisser Mannigfaltigkeiten

 \jour Ber. Verh. S\"achs. Akad. Wiss. Leipzig, Math. Nat. kl.

\yr 1917

\vol 69

\pages 262--277

\elink Available at 

 \url{http://people.csail.mit.edu/bkph/courses/papers/Exact_Conebeam/Radon_Deutsch_1917.pdf}    (February 24, 2016);

%\moreref

%\by  

Radon~J. 

%\paper 

\"Uber die Bestimmung von Funktionen durch ihre Integralwerte l\"angs gewisser Mannigfaltigkeiten,

%\inbook 

{\it Proceedings of Symposia in Applied Mathematics}.

%\bookvol 

vol.~27;

%\ed 

ed.

Lawrence~A. Shepp.

%\publaddr 

Providence,

%\publ 

Amer. Math. Soc.,

%\yr 

1983,

%\pages 

pp.~71--86

\crossref{10.1090/psapm/027/692055}





\Bibitem{dariano:eng}

\by   D'Ariano~G.~M., Mancini~S., Man'ko~V.~I., Tombesi~P.

\paper Reconstructing the density operator by using generalized field quadratures

 \jour Quantum Semiclass. Opt. 

\yr 1996

\issue 5

\vol 8

\pages 1017--1027

\crossref{10.1088/1355-5111/8/5/007}





\Bibitem{leonhardt2:eng}

\by Leonhardt~U.,  Paul~H., D'Ariano~G.~M.

\paper Tomographic reconstruction of the density matrix via pattern functions

 \jour Phys. Rev.~A

\yr 1995

\vol 52

\issue 6

\pages 4899--4907

\crossref{10.1103/physreva.52.4899}



\Bibitem{leonhardt1:eng}

\by  Leonhardt~U., Munroe~M.

\paper Number of phases required to determine a quantum state in optical homodyne tomography

 \jour  Phys. Rev.~A

\yr 1996

\vol 54

\issue 4

\pages 3682--3684

\crossref{10.1103/physreva.54.3682}





\Bibitem{orlow:eng}

\by Or{\l}owski~A.,  Paul~H.

\paper Phase retrieval in quantum mechanics

\jour Phys. Rev. A

\yr 1994

\vol 50

\issue 2

\pages R921--R924

\crossref{10.1103/physreva.50.r921}









\Bibitem{dnn:eng}

\by Amosov~G.~G., Dnestryan~A.~I.

\paper Reconstruction of a~pure state from incomplete information on its optical tomogram

\jour Russian Math. (Iz. VUZ)

\yr 2013

\vol 57

\issue 3

\pages 51--55

\crossref{10.3103/S1066369X13030079}

\scopus{http://www.scopus.com/record/display.url?origin=inward&eid=2-s2.0-84876231479}



\end{thebibliography}



\EngDate

 \renewcommand{\baselinestretch}{0.9}

\newpage

\Russian 

 

 

% 

%\end{document}